\documentclass{article}
\usepackage{enumitem}
\usepackage{etoolbox}
\usepackage{bm}
\usepackage{diagbox}
\usepackage{mathtools}
\usepackage{fancyhdr}
\usepackage{float}
\usepackage{listings}
\lstset{breaklines=true}
\usepackage{graphicx}
\usepackage{xcolor}
\usepackage{titlesec}
\usepackage{titling}
\usepackage{tabularx}
\usepackage{amsmath}
\usepackage{amsfonts}
\usepackage{hyperref}
\usepackage{tikz}
\usepackage[margin=1in]{geometry}

\hypersetup{
	colorlinks=true,
	linkcolor=blue,
	filecolor=magenta,
	urlcolor=blue,
	citecolor=blue,
	anchorcolor=blue
}

\renewcommand{\thesubsubsection}{\Roman{subsubsection}}
\title{SFWRENG-3MX3 Assignment 5}
\author{Matthew Farah, \href{mailto:farahm11@mcmaster.ca}{$\langle$farahm11@mcmaster.ca$\rangle$}, 400468251}
\date{\today}
\makeatletter

\pagestyle{fancy}
\fancyhead{}
\fancyhead[L]{Matthew Farah, $\langle$\href{mailto:farahm11@mcmaster.ca}{farahm11@mcmaster.ca}$\rangle$, 400468251}
\fancyhead[R]{Assignment 5}

\setlength{\parindent}{0cm}

\newcolumntype{Y}{>{\centering\arraybackslash}X}

\newcounter{problem}
\renewcommand{\theproblem}{\arabic{problem}}

\newenvironment{problem}[1][0]{
	\refstepcounter{problem}
	\newpage\vspace{1cm}\large{\par\noindent\textbf{\theproblem) #1}}\par\vspace{.5cm}
	\addcontentsline{toc}{section}{\protect{\large{Problem \theproblem $\;$ \emph{(#1)}}}}
}{\vspace{0.8em}}

\newenvironment{solution}{
	\vspace{0.4cm}\par\noindent\large\textbf{{Solution:}}\par\vspace{1em}
}{\vspace{0.1em}}

\newcounter{partslist}

\makeatletter
\newcommand{\parts@seriesname}{parts@\theproblem}

\newenvironment{parts}{
	\edef\parts@ser{\parts@seriesname}
	\ifcsname\parts@ser\endcsname
		\begin{enumerate}[resume=\parts@ser,label=\textbf{\theproblem.\alph*)},leftmargin=2.5em,itemsep=0.4em,before=\large]
	\else
		\begin{enumerate}[series=\parts@ser,label=\textbf{\theproblem.\alph*)},leftmargin=2.5em,itemsep=0.4em,before=\large]
		\expandafter\gdef\csname\parts@ser\endcsname{1}
	\fi
}{
	\end{enumerate}
}

\makeatother

\makeatletter

\newcounter{currentstep}
\renewcommand{\thecurrentstep}{\Roman{currentstep}}

\newenvironment{steps}{
	\let\steps@olditem\item
	\begin{list}{}{
		\setlength\leftmargin{2.0em}
		\setlength\itemsep{0.4em}
		\setlength\topsep{0.8em}
	}
	\renewcommand{\item}{
		\stepcounter{currentstep}
		\steps@olditem[\textbf{(\thecurrentstep)}]
	}
}{
	\end{list}
	\let\item\steps@olditem
}

\makeatother

\AtBeginEnvironment{parts}{\setcounter{currentstep}{0}}
\AtBeginEnvironment{problem}{\setcounter{currentstep}{0}}

\begin{document}

\maketitle
\pagestyle{fancy}
\tableofcontents

\begin{problem}[Impulse to Frequency Response]
	Use the Discrete Time Fourier Transform (DTFT) or the Continuous Time Fourier Series (CTFT) to find the frequency response that corresponds to the following impulse responses:
\end{problem}

\begin{parts}
	\item $h(t) = \delta(t - 1) + 2\delta(t - 2)$
\end{parts}

\begin{solution}
	\begin{align*}
		X(\omega) &= \int_{-\infty}^{\infty}(x(t) e^{-i\omega{t}}) \, dt \\
		&= \int_{-\infty}^{\infty}(\delta(t - 1) + 2\delta(t - 2)e^{-i\omega{t}}) \, dt \\
		&= \int_{-\infty}^{\infty}(\delta(t - 1)e^{-i\omega{t}} + 2\delta(t - 2)e^{-i\omega{t}}) \, dt \\
		&= e^{-i\omega} + 2e^{-2i\omega} \\
	\end{align*}
\end{solution}

\begin{parts}
	\item $h(t) = \begin{cases} \frac{1}{3}, \; \text{if} \; n = 0, 1, 2 \\ 0, \; \text{otherwise} \end{cases}$
\end{parts}

\begin{solution}

\end{solution}

\begin{problem}[The Step Response]
	In control theory, people use the step response in addition to the impulse response. The reason for this is that the step response gives information on how a controlled system reacts to a change of its set-point, which is a desired state for the system to be in. The step function is given by:

	\begin{align*}
		\text{step}(n) &= \begin{cases} 0, \; \text{if} \; n < 0 \\ 1, \; \text{otherwise} \end{cases}
	\end{align*}

	and the step response is the output of an LTI system in its zero state when $x(n) = \text{step}(n)$.
\end{problem}

\begin{parts}
	\item How might we compute the step response from the impulse response?
\end{parts}

\begin{solution}

\end{solution}

\begin{parts}
	\item Compute the step response of the system given by $y(n) = x(n) + \frac{1}{2}y(n - 1)$ \emph{directly}.
\end{parts}

\begin{solution}

\end{solution}

\begin{parts}
	\item Compute the step response of the system given by $y(n) = x(n) + \frac{1}{2}y(n - 1)$ using the formula derived in \textbf{(a)}.
\end{parts}

\begin{solution}

\end{solution}

\begin{problem}[Filter Design]
	Design a simple FIR filter that removes any frequency component at $\frac{\pi}{2}$ and passes any constant signal with a gain of 1. Give the difference equation and state space representations for such a filter.
\end{problem}

\begin{solution}

\end{solution}

\begin{problem}[Filtering Signals]
	What is the the output of a system with the frequency response:

	\begin{align*}
		H(\omega) &= \begin{cases} 1, \; \text{if} \; |w| \le 2.5 \; \text{radians/second} \end{cases} \\
	\end{align*}

	when its input is a sawtooth signal, the periodic extension of the signal:

	\begin{align*}
		x(t) &= \begin{cases} t, \; \text{if} 0 \le t \le 1 \\ 0, \; \text{otherwise} \end{cases} \\
	\end{align*}

	\emph{Hint}: The sawtooth signal is 1-periodic, hence we can use the Fourier series to determine its frequency content.
\end{problem}

\begin{solution}

\end{solution}

\end{document}
